\section{Methods}
\label{sec:methods}

\subsection{Gene panel selection}
\label{sec:methods_genes}

We selected ten genes spanning five functional categories relevant to spaceflight
radiation biology (\Cref{fig:overview}A; Supplementary Table~S1). Four
\textit{control} genes---BRCA1, TP53, CHEK2, and DNMT3A---were chosen for their
extensive deep mutational scanning (DMS) datasets and well-populated ClinVar
entries, enabling rigorous calibration and benchmarking. Six
\textit{spaceflight-associated} genes---TERT, ATM, NFE2L2, CLOCK, MSTN, and
RAD51---were selected based on their roles in the DNA damage response, telomere
maintenance, oxidative stress signalling, circadian regulation, and muscle
homeostasis, all of which are perturbed during
spaceflight~\cite{Rutter_2024,daSilveira_2020,Twins_2019}. Gene boundaries were
defined as the full genomic locus (5$'$ UTR through 3$'$ UTR plus 2\,kb
flanking sequence on each side) using NCBI RefSeq annotations on the GRCh38
reference genome.

\subsection{Variant generation}
\label{sec:methods_variants}

All possible single-nucleotide variants (SNVs) were enumerated at every
position within each gene locus, yielding three alternate alleles per position.
Indels were generated systematically: single-nucleotide insertions (four bases
at each position) and deletions of length 1--10\,\bp{} at each position within
coding exons, plus frameshifting deletions at selected positions in non-coding
regions. In total, 215,001 variants were generated across the ten genes
(168,409 SNVs and 46,592 indels; Supplementary Table~S2).

\subsection{Evo2 variant effect scoring}
\label{sec:methods_scoring}

Variant effects were predicted using the Evo2 7-billion-parameter genomic
foundation model~\cite{Evo2_2026}, which was trained on the OpenGenome2 dataset
comprising 9.3 trillion nucleotides from 128,000 genomes spanning all domains of
life. The model uses a StripedHyena architecture supporting context windows up to
one million base pairs. For each variant, we extracted the reference and alternate
sequences within a window of $W = 8{,}192$\,\bp{} centred on the variant
position. The model computed per-nucleotide log-likelihoods for both sequences in
a single forward pass. The variant effect score was defined as:

\begin{equation}
  \Delta\text{LL} = \sum_{i=1}^{W}
    \bigl[\log P(x_i^{\text{alt}} \mid \mathbf{x}_{\text{alt}})
    - \log P(x_i^{\text{ref}} \mid \mathbf{x}_{\text{ref}})\bigr]
  \label{eq:delta_ll}
\end{equation}

\noindent where more negative $\Delta$LL values indicate greater predicted
deleteriousness. This zero-shot approach requires no training data and is
agnostic to variant type: the same procedure scores missense, synonymous,
intronic, promoter, and indel variants.

\subsubsection*{Window size optimization.}

We evaluated scoring windows of 1,024, 2,048, 4,096, and 8,192\,\bp{} on
ClinVar-annotated variants in three genes (BRCA1, TP53, CHEK2) using AUROC as
the performance metric. The 8,192\,\bp{} window achieved the highest mean AUROC
(0.9921 across three genes) and was selected for all subsequent analyses
(\Cref{fig:supp_ablation}). Note that 8,192\,\bp{} is the scoring window, not
the model's maximum context; the model can accept inputs up to 1\,Mb but
marginal gains beyond 8,192\,\bp{} did not justify the additional compute cost.

\subsubsection*{Computational resources.}

All inference was performed on NVIDIA A40 (48\,GB VRAM) and A100 (40/80\,GB
VRAM) GPUs on the Cornell Cayuga high-performance computing cluster. Scoring
a single gene required 8--72 hours depending on locus size, totalling
approximately 400 GPU-hours for all ten genes.

\subsection{DMS calibration}
\label{sec:methods_dms}

To calibrate Evo2 scores against experimentally measured variant effects, we
obtained DMS datasets for four control genes from MaveDB~\cite{MaveDB_2019}:

\begin{itemize}
  \item \textbf{BRCA1}: Findlay et al.\ saturation genome editing
    (SGE)~\cite{BRCA1_DMS_2018} (urn:mavedb:00000097-0-2; $n = 3{,}884$
    variants mapped to GRCh38).
  \item \textbf{TP53}: Giacomelli et al.\ nutlin-3 selection
    assay~\cite{TP53_DMS_2018} (urn:mavedb:00000068-b-1; $n = 2{,}827$ matched
    to Evo2-scored positions). In this assay, high functional scores indicate
    loss of function, so the expected correlation with Evo2 $\Delta$LL is
    \textit{negative}.
  \item \textbf{CHEK2}: McCarthy-Leo et al.\ protein abundance
    DMS~\cite{CHEK2_DMS_2024} (urn:mavedb:00001203-a-1; $n = 4{,}882$ matched).
  \item \textbf{DNMT3A}: Garcia et al.\ paired fitness DMS~\cite{DNMT3A_DMS_2025}
    (bioRxiv 10.1101/2025.09.24.678339; $n = 746$ matched). Library-relative
    amino acid positions were converted to absolute DNMT3A coordinates using
    experimentally verified offsets (Library~1: +655, Library~2: +809, Library~3:
    +831); all reference allele matches were confirmed at 100\% against GRCh38.
\end{itemize}

\noindent Spearman rank correlations ($\rho$) between Evo2 $\Delta$LL and DMS
functional scores were computed for each gene. To contextualise Evo2 performance,
we computed the absolute Spearman $|\rho|$ for five additional tools
(AlphaMissense~\cite{AlphaMissense_2023}, CADD~\cite{CADD_2024},
REVEL~\cite{REVEL_2016}, SpliceAI~\cite{SpliceAI_2019}, ncER~\cite{ncER_2019})
on the same variant sets.

\subsection{ClinVar benchmarking}
\label{sec:methods_clinvar}

ClinVar~\cite{ClinVar_2018} variants (release 2024-12) within each gene locus
were filtered to those with $\geq$2-star review status and classified as
pathogenic/likely pathogenic (P/LP) or benign/likely benign (B/LB). Variants of
uncertain significance (VUS) were excluded.

\subsubsection*{Per-tool AUROC.}

For each tool, the area under the receiver operating characteristic curve
(AUROC) was computed using all ClinVar variants scoreable by that tool. This
maximises the sample size available to each tool but means different tools are
evaluated on partially different variant sets.

\subsubsection*{Matched-variant AUROC.}

To enable direct head-to-head comparison, we computed AUROCs on the intersection
of variants scored by all four SNV-capable tools (Evo2 $\cap$ CADD $\cap$
AlphaMissense $\cap$ REVEL). Because AlphaMissense and REVEL are restricted to
missense variants, these matched sets are predominantly missense, which
inherently favours missense-specialised tools.

\subsubsection*{Bootstrap confidence intervals.}

95\% confidence intervals for all AUROCs were computed using stratified bootstrap
resampling (1,000 iterations), maintaining the pathogenic:benign ratio in each
resample.

\subsubsection*{DeLong test.}

Pairwise AUROC differences between tools on matched variant sets were assessed
using the DeLong test~\cite{DeLong_1988}, which accounts for correlation between
ROC curves computed on overlapping data.

\subsection{Orthogonality and ensemble analysis}
\label{sec:methods_ensemble}

Orthogonality between Evo2 and each supervised tool was quantified using Spearman
rank correlation on all shared variants per gene. Ensemble performance was
evaluated using 5-fold cross-validated logistic regression combining Evo2
$\Delta$LL and CADD PHRED scores as features, with ClinVar P/LP vs.\ B/LB as
the binary outcome.

\subsection{Non-coding validation}
\label{sec:methods_noncoding}

\subsubsection*{TERT promoter MPRA.}

Massively parallel reporter assay (MPRA) data for the TERT promoter were
obtained from Kircher et al.~\cite{TERT_MPRA_2019}
(urn:mavedb:00000031-b-1; GBM cell line). Of 973 variants assayed, 777 mapped
to the Evo2-scored TERT locus. Because TERT is transcribed on the minus strand,
all coordinates were maintained on the forward strand with appropriate complement
operations. Spearman correlations between Evo2 $\Delta$LL and MPRA reporter
activity were computed, and the same was done for CADD, SpliceAI, and ncER on
their respective scoreable subsets.

\subsubsection*{ENCODE cCRE enrichment.}

Candidate cis-regulatory elements (cCREs) from the ENCODE
registry~\cite{ENCODE_2020} were intersected with each gene locus. Variants
falling within annotated cCREs---proximal enhancer-like signatures (pELS), distal
enhancer-like signatures (dELS), DNase-H3K4me3, CTCF-bound, and
DNase-only---were compared against non-cCRE variants using mean Evo2 $\Delta$LL,
and enrichment ratios were computed.

\subsection{Mutation signature analysis}
\label{sec:methods_radiation}

Variants were classified into four mutational categories: transitions (Ti),
transversions (Tv), frameshift indels, and in-frame indels. We first tested
whether transversions received more negative (more damaging) Evo2 scores than
transitions using the Mann--Whitney $U$ test.

\subsubsection*{Transversion control.}

To assess whether radiation-signature mutations (C$>$A and G$>$T, associated
with oxidative DNA damage from ionising radiation) score differently from other
transversions of comparable chemical disruption, we performed a
within-transversion comparison. Radiation-oxidative transversions (C$>$A, G$>$T)
were compared against all other transversions (C$>$G, G$>$C, A$>$T, T$>$A,
A$>$C, T$>$G) using the Mann--Whitney $U$ test with rank-biserial correlation as
the effect size. $P$-values were Bonferroni-corrected for ten independent tests.

\subsection{Statistical analysis}
\label{sec:methods_stats}

All statistical tests were two-sided. Spearman rank correlations were used
throughout to avoid distributional assumptions. Mann--Whitney $U$ tests were used
for two-group comparisons of non-normally distributed scores. Bootstrap
confidence intervals used 1,000 stratified resamples. DeLong tests for AUROC
comparison used the fast Sun \& Xu (2014) variance estimator. Multiple testing
was corrected using the Bonferroni method. All analyses were performed in Python
3.11 using NumPy, SciPy, scikit-learn, and custom scripts available at
\url{https://github.com/JangKeunKim/evo2-spaceflight-vep}.

\subsection{Data and code availability}
\label{sec:methods_data}

Pre-computed Evo2 $\Delta$LL scores for all 215,001 variants, ClinVar
benchmarking results, DMS calibration data, and all figure-generation scripts
are deposited at \url{https://github.com/JangKeunKim/evo2-spaceflight-vep}.
The Evo2 model checkpoint is publicly available via Hugging Face
(Arc Institute). ClinVar data are available from NCBI. AlphaMissense, CADD v1.7,
REVEL, SpliceAI, and ncER scores were obtained from their respective public
repositories. DMS and MPRA data are available from
MaveDB~\cite{MaveDB_2019}. ENCODE cCRE annotations are available from
the ENCODE portal~\cite{ENCODE_2020}.
